\section{Graph-Theoretic Analysis}\label{sec:graph-analysis}
\label{sec:reduction-summary}

This section is pure graph theory. We first remove the vertices of large
degree, so that only bounded-degree graphs remain. After that, every case uses
the same idea: a small subgraph that allows a contradiction, that is, a small
subgraph with at least one sign pattern for which its clauses alone are
unsatisfiable. Since such a subgraph has few edges, the random signs hit this
pattern with probability bounded below by a positive constant. We use two kinds
of small subgraphs, dumbbells and graphs of cycle rank at least \(3\), and the
last two subsections show that a graph which survives the reductions contains
linearly many disjoint ones. The random signs then hit the pattern on at least
one of them with probability exponentially close to \(1\).

\subsection{Eliminating vertices of large degree}\label{subsec:bounded-core}

The first reduction turns an arbitrary graph of linear excess into a graph whose
degrees lie in a fixed interval, while the dense case is handled by the
first-moment bound, Lemma~\ref{lem:overview-first-moment}.

\begin{boundedcorerestatement}
Fix \(\excpar>0\). There are constants \(\degb<\infty\) and \(\excd>0\)
depending only on \(\excpar\) such that for every graph \(\Gin\) with
\(\nvert\) vertices, fewer than \(2.41\nvert\) edges and
\(\exc{\Gin}\ge\excpar\nvert\), one can construct a finite multigraph \(H\) with
\[
  |V(H)|\ge\frac{\excpar}{\degb}\,\nvert,
  \qquad
  2\le\dg{H}{v}\le\degb,
  \qquad
  \exc{H}\ge2\excd|V(H)|,
\]
and \(\Prx{\rform{\Gin}\SAT}\le\Prx{\rform{H}\SAT}\).
\end{boundedcorerestatement}

\begin{lemma}[Deleting low-degree vertices]\label{lem:leaf}
Deleting an isolated vertex does not change satisfiability and increases the
excess by \(3/2\). Deleting a vertex of degree \(1\) together with its only
incident edge also does not change satisfiability and increases the excess by
\(1/2\).
\end{lemma}

\begin{proof}
An isolated variable does not occur in the formula. If \(\bvar{v}\) occurs in
only one clause, then every assignment of the other variables can be extended
by choosing the value of \(\bvar{v}\) that makes this clause true.
The two changes in excess are
\[
  \medge-\frac32(\nvert-1)-\left(\medge-\frac32\nvert\right)=\frac32
\]
and
\[
  (\medge-1)-\frac32(\nvert-1)-\left(\medge-\frac32\nvert\right)=\frac12.
\]
\end{proof}

\begin{lemma}[Splitting vertices]\label{lem:splitting}
Fix an integer \(\degb\ge4\). Every finite multigraph \(\Gin\) can be changed
into a finite multigraph \(H\) with maximum degree at most \(\degb\). The graph
\(H\) has the same number of new vertices as new edges:
\[
  |V(H)|-|V(\Gin)|
  =|E(H)|-|E(\Gin)|
  \le\frac{2|E(\Gin)|}{\degb-2}.
\]
Also, there is a coupling of the random signs for which
\[
  \{\rform{\Gin}\SAT\}
  \subseteq
  \{\rform{H}\SAT\}.
\]
Under this coupling, the two signs on every edge of \(H\) are still
independent and uniform, and different edges are independent.
\end{lemma}

\begin{proof}
Recall that a self-loop adds \(2\) to the degree. Write \(b:=\degb-2\ge2\) for
the number of original edge ends that we allow at one copy. For each vertex
\(v\) with \(\dg{\Gin}{v}>b\), let
\[
  k_v=\left\lceil\frac{\dg{\Gin}{v}}{b}\right\rceil,
\]
and replace \(v\) by \(k_v\) copies: split the \(\dg{\Gin}{v}\) edge ends at
\(v\) into groups of size at most \(b\), give one group to each copy, and join
the copies into a path. (The two ends of a self-loop may go to the same copy or
to different copies.) The number of added vertices and the number of added path
edges are both
\[
  N:=\sum_v(k_v-1)
  \le\sum_v\frac{\dg{\Gin}{v}}{b}
  =\frac{2|E(\Gin)|}{\degb-2}.
\]
Each copy has at most \(b\) original edge ends and at most two path edges, so
the maximum degree of \(H\) is at most \(b+2=\degb\).

We now construct the coupling. Fix one vertex \(v\) that was split, and let
\(1,\ldots,k_v\) be its copies along the path. On each path edge \(ij\) we
first put the two clauses
\[
  (\neg\bvar{i}\vee\bvar{j})\wedge(\bvar{i}\vee\neg\bvar{j}),
\]
which force \(\bvar{i}=\bvar{j}\). Together these clauses force all copies to
take the same value, so the new formula is satisfiable exactly when the
original one is. Next, we choose an independent fair bit \(g_i\) for each copy
\(i\) and replace \(\bvar{i}\) by \(\bvar{i}\oplus g_i\) in every clause, which
does not change satisfiability. Finally, on every path edge we keep only one of
its two clauses, chosen independently and uniformly. The result is a
subformula of the previous one, so it is satisfiable whenever the original
formula is.

It remains to check that the resulting signs have the right distribution. Fix
one path of copies and a target sign pattern on all its path edges. If we
choose \(g_i\) at one end of the path and move along the path, then on each
edge the target signs and the bit already chosen determine both the clause that
must be kept and the bit \(g_j\) at the next copy. Thus every target pattern has
exactly two preimages, one for each choice of the first bit, and since all
preimages are equally likely, the signs on the path edges are uniform.
Different paths use independent bits. Finally, the signs on the original edges
were uniform and independent of everything else, and given all bits \(g_i\),
the xor with these bits keeps them uniform and independent of the path edges.
Hence the result has the same distribution as \(\rform{H}\).
\end{proof}

\begin{proof}[Proof of Lemma~\ref{lem:overview-bounded-core}]
Since \(\medge<2.41\,\nvert\), we have
\(\excpar\nvert\le\exc{\Gin}=\medge-\frac32\nvert<0.91\,\nvert\), so
\(\excpar<1\). Fix an integer
\[
  \degb>2+\frac{5}{\excpar} ,
\]
so that \(\degb\ge8\), and apply Lemma~\ref{lem:splitting}; call the result
\(H'\). Let \(N:=|V(H')|-\nvert\) be the number of added vertices. Since it is
also the number of added edges, the excess goes down by \(N/2\), and since
\[
  N\le\frac{2\medge}{\degb-2}<\frac{2\cdot2.41}{5}\,\excpar\nvert<\excpar\nvert,
\]
we get
\[
  \exc{H'}=\exc{\Gin}-\frac N2>\frac{\excpar}{2}\,\nvert,
  \qquad
  |V(H')|=\nvert+N<(1+\excpar)\nvert .
\]
Now we repeatedly apply Lemma~\ref{lem:leaf} to delete vertices of degree \(0\)
or \(1\), and call the resulting graph \(H\). Since this step can only increase
the excess and cannot increase any degree or the number of vertices, we get
\[
  2\le\dg{H}{v}\le\degb,
  \qquad
  \exc{H}\ge\frac{\excpar}{2}\,\nvert
  >\frac{\excpar}{2(1+\excpar)}\,|V(H)| .
\]
So we may take
\[
  \excd:=\frac{\excpar}{4(1+\excpar)} .
\]
The degree bound gives
\[
  \frac{\excpar}{2}\,\nvert
  \le\exc{H}
  \le |E(H)|
  \le\frac{\degb}{2}|V(H)|,
\]
so \(|V(H)|\ge(\excpar/\degb)\nvert\). Finally,
Lemmas~\ref{lem:splitting} and~\ref{lem:leaf} give
\(\Prx{\rform{\Gin}\SAT}\le\Prx{\rform{H}\SAT}\), and both \(\degb\) and
\(\excd\) depend only on \(\excpar\).
\end{proof}

\subsection{Unsatisfiable support}\label{subsec:unsat-support}

We say that a finite multigraph \(H\) supports an unsatisfiable 2-CNF if there
is at least one sign pattern on the edges of \(H\) for which the resulting
formula is unsatisfiable. This is a property of the graph alone. Every sign
pattern on \(H\) has positive probability, so it is the same as saying that
\(\rform{H}\) is unsatisfiable with positive probability.

We use two sufficient conditions, which are the two cases below. The first is a
\emph{dumbbell}, that is, a self-loop joined by a path to a second cycle. With
suitable signs each part has one job: the self-loop fixes its variable, the
path carries the forced value, and the second cycle gives the contradiction.
Without a self-loop we ask for cycle rank at least \(3\), because cycle rank
\(2\) is not always enough there: three parallel edges between two vertices
have cycle rank \(2\), but they are satisfiable under every sign pattern.

\begin{unsatsupportrestatement}
Let \(\Gin\) be a finite connected multigraph. Then \(\Gin\) supports an
unsatisfiable 2-CNF, that is, there is a sign pattern on its edges for which
the resulting formula is unsatisfiable, in either of the following cases:
\begin{enumerate}[label=(\arabic*),leftmargin=3em]
\item \(\Gin\) has no self-loops and \(\cyc{\Gin}\ge3\);
\item \(\Gin\) has a self-loop and \(\cyc{\Gin}\ge2\).
\end{enumerate}
\end{unsatsupportrestatement}

\begin{lemma}[Dumbbells: multigraphs with a self-loop]\label{lem:selfloop-support}
Let \(H\) be a finite connected multigraph with a self-loop. If
\[
  \cyc{H}\ge2,
\]
then \(H\) supports an unsatisfiable 2-CNF.
\end{lemma}

\begin{proof}
Choose a self-loop \(e_0\) at a vertex \(u\), and put \((\bvar{u}\vee\bvar{u})\)
on it, which forces \(\bvar{u}=1\). Since \(\cyc{H}\ge2\), the graph still
contains a cycle after \(e_0\) is deleted, and together with a path from \(u\)
to this cycle, the self-loop and the cycle form a dumbbell. Walk along the path
and then around the cycle. On each edge \(ab\) of this walk except the last
one, put \((\neg\bvar{a}\vee\bvar{b})\), where \(a\) is the vertex reached
first; these clauses carry the value \(1\) along the path and around the cycle.
On the last edge of the cycle put \((\neg\bvar{a}\vee\neg\bvar{b})\), which
gives a contradiction. (If the cycle is itself a self-loop at \(a\), put
\((\neg\bvar{a}\vee\neg\bvar{a})\) on it.) Arbitrary clauses on the unused
edges complete the formula.
\end{proof}

\begin{lemma}[Multigraphs without self-loops]\label{lem:multisupport}
Let \(H\) be a finite connected multigraph without self-loops. If
\[
  \cyc{H}\ge3,
\]
then \(H\) supports an unsatisfiable 2-CNF.
\end{lemma}

\begin{proof}
For each edge \(e=uv\) of \(H\), insert a new vertex \(y_e\) used only for that
edge, and replace \(e\) by the path \(u-y_e-v\). Since parallel edges get
different new vertices, the resulting graph \(\widetilde H\) is simple, and
since each subdivision adds one vertex and one edge,
\[
  \cyc{\widetilde H}=\cyc{H}\ge3.
\]
By Karve--Hirani~\cite[Corollary 19]{KarveHirani2019},
\(\widetilde H\) supports an unsatisfiable formula \(\widetilde\fml\).
On the path \(u-y_e-v\) this formula has two clauses: one contains a literal
\(a_e\) of \(\bvar{u}\) and a literal of \(\bvar{y_e}\), and the other contains
a literal \(b_e\) of \(\bvar{v}\) and a literal of \(\bvar{y_e}\). Since
\(\bvar{y_e}\) appears only in these two clauses, we can eliminate it. If
\(\bvar{y_e}\) has opposite signs in the two clauses, then
\[
  \exists \bvar{y_e}\,
  \bigl[(a_e\vee \bvar{y_e})\wedge(b_e\vee \neg\bvar{y_e})\bigr]
  \equiv
  a_e\vee b_e ,
\]
and the case with the two signs exchanged is the same. If \(\bvar{y_e}\) has the
same sign in both clauses, then one value of \(\bvar{y_e}\) satisfies both
clauses, and the pair can be replaced by the always-true formula \(\top\).
After all new variables are eliminated, \(\widetilde\fml\) becomes an
unsatisfiable formula that uses each edge of \(H\) at most once, and adding
arbitrary clauses on the unused edges keeps it unsatisfiable. Hence \(H\)
supports an unsatisfiable formula.
\end{proof}

The two lemmas above prove Lemma~\ref{lem:overview-unsat-support}. We now prove
the two statements that use it, namely the positive-excess component reduction
and the packing statement. Recall from Section~\ref{subsec:framework} that when
the maximum degree is at most \(\degb\), every ball of radius \(\rad\) has at
most \(\ballsz{\rad}\) vertices.

\begin{poscomponentsrestatement}
Let \(\Gin\) be a finite multigraph with maximum degree at most \(\degb\).
Delete every connected component \(C\) with \(\exc{C}\le0\), and call the
result \(H\). Then every vertex of \(H\) has the same degree in \(H\) as in
\(\Gin\),
\[
  \exc{H}\ge\exc{\Gin},
  \qquad
  |V(H)|\ge\frac{2}{\degb}\,\exc{\Gin},
\]
every component of \(H\) without a self-loop has cycle rank at least \(3\),
every component with a self-loop has cycle rank at least \(2\), and
\(\Prx{\rform{\Gin}\SAT}\le\Prx{\rform{H}\SAT}\).
\end{poscomponentsrestatement}

\begin{proof}
Since the excess of a graph is the sum of the excesses of its components,
deleting components with nonpositive excess gives
\[
  \exc{H}
  =\exc{\Gin}-\sum_{C:\,\exc{C}\le0}\exc{C}
  \ge\exc{\Gin}.
\]
Because we delete whole components, the degrees of the remaining vertices do
not change, and because we only delete clauses,
\(\{\rform{\Gin}\SAT\}\subseteq\{\rform{H}\SAT\}\). The degree bound gives
\[
  \exc{\Gin}
  \le\exc{H}
  \le|E(H)|
  \le\frac{\degb}{2}|V(H)| .
\]

For every connected component \(C\) of \(H\),
\[
  0<\exc{C}
  =\cyc{C}-1-\frac{|V(C)|}{2}.
\]
If \(C\) has no self-loops, then \(|V(C)|\ge2\), because a single vertex
without edges has negative excess; hence \(\cyc{C}>2\), and since \(\cyc{C}\)
is an integer, \(\cyc{C}\ge3\). If \(C\) has a self-loop, then \(|V(C)|\ge1\),
so \(\cyc{C}>3/2\) and \(\cyc{C}\ge2\).
\end{proof}

\begin{packingrestatement}
Fix \(\degb<\infty\), an integer \(\radp\ge1\) and \(\pack>0\). Suppose \(\Gin\)
has \(\nvert\) vertices and maximum degree at most \(\degb\), and that at least
\(\pack\nvert\) vertices \(v\) satisfy one of the following:
\begin{enumerate}[label=(\arabic*),leftmargin=3em]
\item \(\ball{\radp}{v}\) has no self-loop and \(\cyc{\ball{\radp}{v}}\ge3\);
\item \(\ball{\radp}{v}\) has a self-loop and \(\cyc{\ball{\radp}{v}}\ge2\).
\end{enumerate}
Then, with the implied constant depending only on \(\degb,\radp,\pack\),
\[
  \Prx{\rform{\Gin}\SAT}\le e^{-\Omega(\nvert)} .
\]
\end{packingrestatement}

\begin{proof}
Call the vertices in the hypothesis candidates, and choose centers greedily
among them: each time a center is chosen, remove every candidate at distance at
most \(2\radp\) from it. At most \(\ballsz{2\radp}\) candidates are removed
each time, and two chosen centers are at distance more than \(2\radp\), so
their balls of radius \(\radp\) are vertex-disjoint. Thus we obtain at least
\[
  \frac{\pack \nvert}{\ballsz{2\radp}}
\]
pairwise vertex-disjoint balls of radius \(\radp\).

Each ball is connected, because a shortest path from the center to any vertex of
the ball stays inside the ball. Hence, by
Lemma~\ref{lem:overview-unsat-support}, every ball supports an unsatisfiable
2-CNF, and we fix one such sign pattern for each ball. A ball has at most
\(\degb\ballsz{\radp}/2\) edges, and each edge has four sign choices, so the
random signs on the ball equal this pattern with probability at least
\(4^{-\degb\ballsz{\radp}/2}=2^{-\degb\ballsz{\radp}}\). Disjoint balls have
disjoint edge sets, so these events are independent, and if one of them
happens, the formula is unsatisfiable. Therefore,
\[
  \Prx{\rform{\Gin}\SAT}
  \le\bigl(1-2^{-\degb\ballsz{\radp}}\bigr)^{\pack \nvert/\ballsz{2\radp}}
  \le \exp\Bigl(-2^{-\degb\ballsz{\radp}}\,\frac{\pack \nvert}{\ballsz{2\radp}}\Bigr).
\]
\end{proof}

\subsection{Self-loop reduction}\label{subsec:loop-reduction}

\begin{loopcountingrestatement}
Fix \(\degb<\infty\) and \(\excd>0\). There are \(\radp<\infty\) and
\(\pack>0\) depending only on them with the following property. Let \(\Gin\) be
an \(\nvert\)-vertex multigraph with maximum degree at most \(\degb\) in which
every component containing a self-loop has cycle rank at least \(2\). If \(\Gin\)
has more than \(\excd\nvert\) self-loops, then at least \(\pack\nvert\) vertices
\(v\) have a ball \(\ball{\radp}{v}\) that contains a self-loop and has cycle
rank at least \(2\).
\end{loopcountingrestatement}

The proof only counts vertices and distances.

\begin{proof}[Proof of Lemma~\ref{lem:overview-loop-counting}]
Let \(V_{\mathrm{loop}}\) be the set of vertices that have at least one
self-loop. Since each vertex has at most \(\degb/2\) self-loops,
\[
  |V_{\mathrm{loop}}|>\frac{2\excd}{\degb}\,\nvert .
\]
Choose an integer \(\radp\) large enough that
\[
  \frac1{\lfloor \radp/2\rfloor+1}\le\frac{\excd}{\degb}.
\]

Call \(v\in V_{\mathrm{loop}}\) good if at least one of the following holds:
\begin{enumerate}[label=(\arabic*),leftmargin=3em]
\item \(v\) has at least two self-loops;
\item the ball \(\ball{\radp}{v}\) contains the whole component of \(v\);
\item another vertex of \(V_{\mathrm{loop}}\) is at distance at most \(\radp\) from \(v\).
\end{enumerate}
For every good vertex \(v\), the ball \(\ball{\radp}{v}\) contains a self-loop
and has cycle rank at least \(2\): in cases~(1) and~(3) the ball contains two
self-loops, and in case~(2) it is the whole component of \(v\), which has
cycle rank at least \(2\) by assumption.

We now bound the number of bad vertices. Two bad vertices are at distance more
than \(\radp\), so the balls of radius
\[
  R':=\left\lfloor\frac{\radp}2\right\rfloor
\]
around them are pairwise vertex-disjoint. Each such ball has at least \(R'+1\)
vertices: since \(\ball{\radp}{v}\) does not contain the whole component of
\(v\), there is a shortest path from \(v\) to a vertex at distance more than
\(\radp\), and its first \(R'+1\) vertices lie in \(\ball{R'}{v}\). Hence there
are at most \(\nvert/(R'+1)\le(\excd/\degb)\nvert\) bad vertices, and more than
\[
  \frac{2\excd}{\degb}\,\nvert-\frac{\excd}{\degb}\,\nvert
  =\frac{\excd}{\degb}\,\nvert
\]
good vertices. This proves the lemma with \(\pack:=\excd/\degb\).
\end{proof}

\begin{lemma}[Removing self-loops]\label{lem:remove-loops}
Suppose that \(\Gin\) has \(\nvert\) vertices and satisfies
\[
  2\le\dg{\Gin}{v}\le\degb,
  \qquad
  \exc{\Gin}\ge2\excd \nvert,
\]
and that every connected component containing a self-loop has cycle rank at
least \(2\). Then either
\[
  \Prx{\rform{\Gin}\SAT}\le e^{-\Omega(\nvert)},
\]
with the implied constant depending only on \(\degb,\excd\), or one can obtain
a graph \(H\) without self-loops by deleting only clauses and low-degree
vertices, such that
\[
  |V(H)|\ge\frac{2\excd}{\degb}\,\nvert,
  \qquad
  2\le\dg{H}{v}\le\degb,
  \qquad
  \exc{H}\ge\excd|V(H)|,
\]
every connected component of \(H\) has cycle rank at least \(3\), and
\[
  \Prx{\rform{\Gin}\SAT}
  \le
  \Prx{\rform{H}\SAT}.
\]
\end{lemma}

\begin{proof}
First suppose that \(\Gin\) has more than \(\excd\nvert\) self-loops. Then
Lemma~\ref{lem:overview-loop-counting} gives at least \(\pack\nvert\) vertices
whose ball of radius \(\radp\) contains a self-loop and has cycle rank at least
\(2\), and Corollary~\ref{cor:overview-packing} gives the first conclusion.

Now suppose that \(\Gin\) has at most \(\excd\nvert\) self-loops, and delete all
of them. Each deleted self-loop lowers the excess by \(1\), so the excess is
still at least \(\excd\nvert\). Next we repeatedly delete vertices of degree
\(0\) or \(1\) and call the resulting graph \(H'\); by Lemma~\ref{lem:leaf},
this step does not lower the excess, so \(\exc{H'}\ge\excd\nvert\). Finally, we
apply Lemma~\ref{lem:overview-positive-components} to \(H'\) and call the
result \(H\). Then \(H\) has no self-loops, its degrees are the same as in
\(H'\) and hence lie in \([2,\degb]\), every component has cycle rank at least
\(3\), and
\[
  \exc{H}\ge\exc{H'}\ge\excd\nvert\ge\excd|V(H)|,
  \qquad
  |V(H)|\ge\frac{2}{\degb}\,\exc{H'}\ge\frac{2\excd}{\degb}\,\nvert .
\]
Every operation above deletes clauses or keeps satisfiability unchanged, so
\(\Prx{\rform{\Gin}\SAT}\le\Prx{\rform{H}\SAT}\).
\end{proof}

\subsection{Local-tree reduction}\label{subsec:local-tree-reduction}

\begin{finitecountingrestatement}
Fix \(\degb<\infty\), an integer \(\rad\ge1\) and \(\nontree>0\). There are
\(\radp<\infty\) and \(\pack>0\) depending only on them with the following
property. Let \(\Gin\) be an \(\nvert\)-vertex graph without self-loops, with
maximum degree at most \(\degb\), in which every component has cycle rank at
least \(3\). If \(\Gin\) is not \loctree{\rad}{\nontree}, then at least
\(\pack\nvert\) vertices \(v\) have a ball \(\ball{\radp}{v}\) with cycle rank
at least \(3\).
\end{finitecountingrestatement}

\begin{proof}[Proof of Lemma~\ref{lem:overview-finite-counting}]
Since \(\Gin\) is not \loctree{\rad}{\nontree}, more than \(\nontree \nvert\)
vertices \(v\) have a ball \(\ball{\rad}{v}\) that is not a tree. We choose
centers greedily among these vertices, removing every candidate at distance at
most \(2\rad\) each time a center is chosen. This gives
\[
  k\ge\frac{\nontree}{\ballsz{2\rad}}\,\nvert=:a\nvert
\]
centers whose balls of radius \(\rad\) are pairwise vertex-disjoint. None of
these balls is a tree, so we can fix one cycle in each of them, and these \(k\)
cycles are pairwise vertex-disjoint.

Fix an integer \(s\), to be chosen at the end, and build an auxiliary graph on
the \(k\) centers in which two centers are joined if their distance in \(\Gin\)
is at most \(s\). We call the connected components of the auxiliary graph
\emph{cycle groups}, and we set
\[
  \radp:=2s+\rad .
\]

\emph{Case 1: at least \(k/2\) centers lie in cycle groups of size at least
\(3\).} Take such a center \(w\). In the auxiliary graph, two other centers can
be reached from \(w\) in at most two steps, so these two centers are at
distance at most \(2s\) from \(w\). Since their cycles lie within distance
\(\rad\) of them, the ball \(\ball{\radp}{w}\) contains three of the
vertex-disjoint cycles, and its cycle rank is at least \(3\). This gives at
least \(k/2\) good centers.

\emph{Case 2: at least \(k/2\) centers lie in cycle groups of size \(1\) or
\(2\).} Then there are at least \(k/4\) such small groups. We choose one center
\(w\) from each small group and let
\[
  R':=\left\lfloor\frac s2\right\rfloor.
\]
Centers from different groups are at distance more than \(s\), so their balls
of radius \(R'\) are pairwise vertex-disjoint. Suppose that for more than half
of the chosen centers \(w\), the ball \(\ball{R'}{w}\) does not contain the
whole component of \(w\). Then each of these balls has at least \(R'+1\)
vertices, because a shortest path from \(w\) to a vertex outside the ball has
its first \(R'+1\) vertices inside the ball, and since these balls are
disjoint,
\[
  \nvert\ge\frac{k}{8}(R'+1)
  \ge\frac{a\nvert}{8}(R'+1).
\]
We choose \(s\) so large that \(a(R'+1)>8\), which makes this impossible.
Hence for at least half of the chosen centers \(w\), that is, for at least
\(k/8\) centers, the ball \(\ball{R'}{w}\) contains the whole component of
\(w\). Since \(R'\le\radp\), the ball \(\ball{\radp}{w}\) is this component,
and its cycle rank is at least \(3\) by assumption.

In both cases at least \(k/8\ge(a/8)\nvert\) vertices \(w\) satisfy
\(\cyc{\ball{\radp}{w}}\ge3\), which proves the lemma with \(\pack:=a/8\). The
constants \(s\), \(\radp\) and \(\pack\) depend only on \(\degb,\rad,\nontree\).
\end{proof}
